% !TEX root = ../main.tex
\chapter*{Abstract}
\addcontentsline{toc}{chapter}{Abstract}

Global reanalyses describe the atmosphere at a resolution of tens of kilometres, at which the
terrain that decides Swiss weather has been averaged away. This work studies statistical
downscaling of daily maximum temperature and daily precipitation from ERA5 and ERA5-Land to the
one-kilometre MeteoSwiss analysis grids with Convolutional Conditional Neural Processes, an
architecture that predicts a full probability distribution at arbitrary target coordinates. A
Gaussian output head serves temperature and a Bernoulli--Gamma head precipitation. Twelve model
configurations, seven for temperature and five for precipitation, are trained on 2020--2023 under a temporal five-fold cross-validation and
evaluated twice: on their held-out folds and on 2024, a year no fold ever saw. A permutation-based
feature-importance study measures the reliance on the different input channels. Finally, all selected
models are verified against station observations that entered no training or selection step.

For temperature, input built from the atmospheric state, pressure-level temperature, humidity,
geopotential and wind, beats input built from the coarse surface field on every metric, and
generalises better: the atmospheric configurations hold or improve on the unseen year while the surface
configurations degrade, with the largest gap in the cold season, when the fine-scale truth
decouples from the coarse field. The selected model reaches a skill of $0.71$ against bilinear
interpolation and a mean absolute error of $1.03\,^{\circ}\mathrm{C}$ on 2024, and its
importances sit on lower-tropospheric temperature at physically sensible levels and hours rather
than on coordinates or terrain, so the learnt mapping is meteorological rather than a trained
interpolation. For precipitation the value of the models is distributional: the Bernoulli--Gamma
head removes the drizzle bias of the coarse field and restores the observed wet-day frequency and
intensity, worth a ranked probability skill near $+0.4$ over five intensity categories in both
regimes, while in plain millimetre error the reference interpolation remains unbeaten on the
gridded target. Verification at $147$ temperature stations and $139$ rain gauges confirms both
verdicts at real measurement points, where one of the two selected precipitation runs even takes the lower
millimetre error. The heaviest precipitation events and the warm temperature tail remain beyond the reach
of every configuration tested, and marginal per-point distributions leave the spatial coherence
of the fields an open question.
